\iftestbuild % TEST-ONLY-BEGIN
\setcounter{section}{5}\setcounter{theorem}{107}
\fi % TEST-ONLY-END
% Source: book pp. 105--118 (PDF 119--132); solutions pp. 113--128, 172--173.
% A short note set in small italics under an exercise, as Axler does.
\section{Duality}

\subsection{Dual Space and Dual Map}

Linear maps into the scalar field $\F$ play a special role in linear algebra, and
thus they get a special name.

\begin{definition}[Linear functional]\label{ax:3.108}
A \emph{linear functional} on $V$ is a linear map from $V$ to $\F$. In other
words, a linear functional is an element of $\Lin(V,\F)$.
\end{definition}

\begin{example}[Linear functionals]\label{ax:3.109}\leavevmode
\begin{itemize}
\item Define $\varphi\colon\R^3\to\R$ by $\varphi(x,y,z)=4x-5y+2z$. Then
  $\varphi$ is a linear functional on $\R^3$.
\item Fix $(c_1,\dots,c_n)\in\F^n$. Define $\varphi\colon\F^n\to\F$ by
  $\varphi(x_1,\dots,x_n)=c_1x_1+\cdots+c_nx_n$. Then $\varphi$ is a linear
  functional on $\F^n$.
\item Define $\varphi\colon\Poly(\R)\to\R$ by
  \[ \varphi(p) = 3p''(5)+7p(4). \]
  Then $\varphi$ is a linear functional on $\Poly(\R)$.
\item Define $\varphi\colon\Poly(\R)\to\R$ by
  \[ \varphi(p) = \int_0^1 p \]
  for each $p\in\Poly(\R)$. Then $\varphi$ is a linear functional on
  $\Poly(\R)$.
\end{itemize}
\end{example}

The vector space $\Lin(V,\F)$ also gets a special name and special notation.

\begin{definition}[Dual space, $V'$]\label{ax:3.110}
The \emph{dual space} of $V$, denoted by $V'$, is the vector space of all linear
functionals on $V$. In other words, $V'=\Lin(V,\F)$.
\end{definition}

\begin{result}[$\dim V'=\dim V$]\label{ax:3.111}
Suppose $V$ is finite-dimensional. Then $V'$ is also finite-dimensional and
\[ \dim V' = \dim V. \]
\end{result}

\begin{proof}
By \ref{ax:3.72} we have
\[ \dim V' = \dim\Lin(V,\F) = (\dim V)(\dim\F) = \dim V, \]
as desired.
\end{proof}

In the following definition, the linear map lemma (\ref{ax:3.4}) implies that
each $\varphi_j$ is well defined.

\begin{definition}[Dual basis]\label{ax:3.112}
If $v_1,\dots,v_n$ is a basis of $V$, then the \emph{dual basis} of
$v_1,\dots,v_n$ is the list $\varphi_1,\dots,\varphi_n$ of elements of $V'$,
where each $\varphi_j$ is the linear functional on $V$ such that
\[ \varphi_j(v_k) = \begin{cases} 1 & \text{if } k=j,\\
                                 0 & \text{if } k\ne j. \end{cases} \]
\end{definition}

\begin{example}[The dual basis of the standard basis of $\F^n$]\label{ax:3.113}
Suppose $n$ is a positive integer. For $1\le j\le n$, define $\varphi_j$ to be
the linear functional on $\F^n$ that selects the $j^{\text{th}}$ coordinate of a
vector in $\F^n$. Thus
\[ \varphi_j(x_1,\dots,x_n) = x_j \]
for each $(x_1,\dots,x_n)\in\F^n$.

Let $e_1,\dots,e_n$ be the standard basis of $\F^n$. Then
\[ \varphi_j(e_k) = \begin{cases} 1 & \text{if } k=j,\\
                                 0 & \text{if } k\ne j. \end{cases} \]
Thus $\varphi_1,\dots,\varphi_n$ is the dual basis of the standard basis
$e_1,\dots,e_n$ of $\F^n$.
\end{example}

The next result shows that the dual basis of a basis of $V$ consists of the
linear functionals on $V$ that give the coefficients for expressing a vector in
$V$ as a linear combination of the basis vectors.

\begin{result}[Dual basis gives coefficients for linear combination]\label{ax:3.114}
Suppose $v_1,\dots,v_n$ is a basis of $V$ and $\varphi_1,\dots,\varphi_n$ is the
dual basis. Then
\[ v = \varphi_1(v)v_1+\cdots+\varphi_n(v)v_n \]
for each $v\in V$.
\end{result}

\begin{proof}
Suppose $v\in V$. Then there exist $c_1,\dots,c_n\in\F$ such
that\refstepcounter{theorem}
\begin{equation*}
v = c_1v_1+\cdots+c_nv_n. \tag*{\thetheorem}\label{ax:3.115}
\end{equation*}
If $j\in\{1,\dots,n\}$, then applying $\varphi_j$ to both sides of the equation
above gives
\[ \varphi_j(v) = c_j. \]
Substituting the values for $c_1,\dots,c_n$ given by the equation above into
\ref{ax:3.115} shows that $v=\varphi_1(v)v_1+\cdots+\varphi_n(v)v_n$.
\end{proof}

The next result shows that the dual basis is indeed a basis of the dual space.
Thus the terminology ``dual basis'' is justified.

\begin{result}[Dual basis is a basis of the dual space]\label{ax:3.116}
Suppose $V$ is finite-dimensional. Then the dual basis of a basis of $V$ is a
basis of $V'$.
\end{result}

\begin{proof}
Suppose $v_1,\dots,v_n$ is a basis of $V$. Let $\varphi_1,\dots,\varphi_n$ denote
the dual basis.

To show that $\varphi_1,\dots,\varphi_n$ is a linearly independent list of
elements of $V'$, suppose $a_1,\dots,a_n\in\F$ are such
that\refstepcounter{theorem}
\begin{equation*}
a_1\varphi_1+\cdots+a_n\varphi_n = 0. \tag*{\thetheorem}\label{ax:3.117}
\end{equation*}
Now
\[ (a_1\varphi_1+\cdots+a_n\varphi_n)(v_k) = a_k \]
for each $k=1,\dots,n$. Thus \ref{ax:3.117} shows that $a_1=\cdots=a_n=0$. Hence
$\varphi_1,\dots,\varphi_n$ is linearly independent.

Because $\varphi_1,\dots,\varphi_n$ is a linearly independent list in $V'$ whose
length equals $\dim V'$ (by \ref{ax:3.111}), we can conclude that
$\varphi_1,\dots,\varphi_n$ is a basis of $V'$ (see \ref{ax:2.38}).
\end{proof}

In the definition below, note that if $T$ is a linear map from $V$ to $W$ then
$T'$ is a linear map from $W'$ to $V'$.

\begin{definition}[Dual map, $T'$]\label{ax:3.118}
Suppose $T\in\Lin(V,W)$. The \emph{dual map} of $T$ is the linear map
$T'\in\Lin(W',V')$ defined for each $\varphi\in W'$ by
\[ T'(\varphi) = \varphi\circ T. \]
\end{definition}

If $T\in\Lin(V,W)$ and $\varphi\in W'$, then $T'(\varphi)$ is defined above to be
the composition of the linear maps $\varphi$ and $T$. Thus $T'(\varphi)$ is
indeed a linear map from $V$ to $\F$; in other words, $T'(\varphi)\in V'$.

The following two bullet points show that $T'$ is a linear map from $W'$ to
$V'$.
\begin{itemize}
\item If $\varphi,\psi\in W'$, then
  \[ T'(\varphi+\psi) = (\varphi+\psi)\circ T = \varphi\circ T+\psi\circ T
     = T'(\varphi)+T'(\psi). \]
\item If $\lambda\in\F$ and $\varphi\in W'$, then
  \[ T'(\lambda\varphi) = (\lambda\varphi)\circ T = \lambda(\varphi\circ T)
     = \lambda T'(\varphi). \]
\end{itemize}

The prime notation appears with two unrelated meanings in the next example: $D'$
denotes the dual of the linear map $D$, and $p'$ denotes the derivative of a
polynomial $p$.

\begin{example}[Dual map of the differentiation linear map]\label{ax:3.119}
Define $D\colon\Poly(\R)\to\Poly(\R)$ by $Dp=p'$.
\begin{itemize}
\item Suppose $\varphi$ is the linear functional on $\Poly(\R)$ defined by
  $\varphi(p)=p(3)$. Then $D'(\varphi)$ is the linear functional on $\Poly(\R)$
  given by
  \[ \bigl(D'(\varphi)\bigr)(p) = (\varphi\circ D)(p) = \varphi(Dp)
     = \varphi(p') = p'(3). \]
  Thus $D'(\varphi)$ is the linear functional on $\Poly(\R)$ taking $p$ to
  $p'(3)$.
\item Suppose $\varphi$ is the linear functional on $\Poly(\R)$ defined by
  $\varphi(p)=\int_0^1p$. Then $D'(\varphi)$ is the linear functional on
  $\Poly(\R)$ given by
  \begin{align*}
  \bigl(D'(\varphi)\bigr)(p) &= (\varphi\circ D)(p)\\
    &= \varphi(Dp)\\
    &= \varphi(p')\\
    &= \int_0^1 p'\\
    &= p(1)-p(0).
  \end{align*}
  Thus $D'(\varphi)$ is the linear functional on $\Poly(\R)$ taking $p$ to
  $p(1)-p(0)$.
\end{itemize}
\end{example}

In the next result, (a) and (b) imply that the function that takes $T$ to $T'$
is a linear map from $\Lin(V,W)$ to $\Lin(W',V')$.

In (c) below, note the reversal of order from $ST$ on the left to $T'S'$ on the
right.

\begin{result}[Algebraic properties of dual maps]\label{ax:3.120}
Suppose $T\in\Lin(V,W)$. Then
\begin{parts}
\item $(S+T)'=S'+T'$ for all $S\in\Lin(V,W)$;
\item $(\lambda T)'=\lambda T'$ for all $\lambda\in\F$;
\item $(ST)'=T'S'$ for all $S\in\Lin(W,U)$.
\end{parts}
\end{result}

\begin{proof}
The proofs of (a) and (b) are left to the reader.

To prove (c), suppose $\varphi\in U'$. Then
\[ (ST)'(\varphi) = \varphi\circ(ST) = (\varphi\circ S)\circ T
   = T'(\varphi\circ S) = T'\bigl(S'(\varphi)\bigr) = (T'S')(\varphi), \]
\marginnote{Some books use the notation $V^*$ and $T^*$ for duality instead of
$V'$ and $T'$. However, here we reserve the notation $T^*$ for the adjoint,
which will be introduced when we study linear maps on inner product spaces in
Chapter~7.}%
where the first, third, and fourth equalities above hold because of the
definition of the dual map, the second equality holds because composition of
functions is associative, and the last equality follows from the definition of
composition.

The equation above shows that $(ST)'(\varphi)=(T'S')(\varphi)$ for all
$\varphi\in U'$. Thus $(ST)'=T'S'$.
\end{proof}

\subsection{Null Space and Range of Dual of Linear Map}

Our goal in this subsection is to describe $\nul T'$ and $\range T'$ in terms of
$\range T$ and $\nul T$. To do this, we will need the next definition.

\begin{definition}[Annihilator, $U^0$]\label{ax:3.121}
For $U\subseteq V$, the \emph{annihilator} of $U$, denoted by $U^0$, is defined
by
\[ U^0 = \{\varphi\in V' : \varphi(u)=0 \text{ for all } u\in U\}. \]
\end{definition}

\begin{example}[Element of an annihilator]\label{ax:3.122}
Suppose $U$ is the subspace of $\Poly(\R)$ consisting of polynomial multiples of
$x^2$. If $\varphi$ is the linear functional on $\Poly(\R)$ defined by
$\varphi(p)=p'(0)$, then $\varphi\in U^0$.
\end{example}

For $U\subseteq V$, the annihilator $U^0$ is a subset of the dual space $V'$.
Thus $U^0$ depends on the vector space containing $U$, so a notation such as
$U_V^0$ would be more precise. However, the containing vector space will always
be clear from the context, so we will use the simpler notation $U^0$.

\begin{example}[The annihilator of a two-dimensional subspace of
$\R^5$]\label{ax:3.123}
Let $e_1,e_2,e_3,e_4,e_5$ denote the standard basis of $\R^5$; let
$\varphi_1,\varphi_2,\varphi_3,\varphi_4,\varphi_5\in(\R^5)'$ denote the dual
basis of $e_1,e_2,e_3,e_4,e_5$. Suppose
\[ U = \spn(e_1,e_2) = \{(x_1,x_2,0,0,0)\in\R^5 : x_1,x_2\in\R\}. \]
We want to show that $U^0=\spn(\varphi_3,\varphi_4,\varphi_5)$.

Recall (see \ref{ax:3.113}) that $\varphi_j$ is the linear functional on $\R^5$
that selects the $j^{\text{th}}$ coordinate:
$\varphi_j(x_1,x_2,x_3,x_4,x_5)=x_j$.

First suppose $\varphi\in\spn(\varphi_3,\varphi_4,\varphi_5)$. Then there exist
$c_3,c_4,c_5\in\R$ such that $\varphi=c_3\varphi_3+c_4\varphi_4+c_5\varphi_5$.
If $(x_1,x_2,0,0,0)\in U$, then
\[ \varphi(x_1,x_2,0,0,0)
   = (c_3\varphi_3+c_4\varphi_4+c_5\varphi_5)(x_1,x_2,0,0,0) = 0. \]
Thus $\varphi\in U^0$. Hence we have shown that
$\spn(\varphi_3,\varphi_4,\varphi_5)\subseteq U^0$.

To show the inclusion in the other direction, suppose that $\varphi\in U^0$.
Because the dual basis is a basis of $(\R^5)'$, there exist
$c_1,c_2,c_3,c_4,c_5\in\R$ such that
$\varphi=c_1\varphi_1+c_2\varphi_2+c_3\varphi_3+c_4\varphi_4+c_5\varphi_5$.
Because $e_1\in U$ and $\varphi\in U^0$, we have
\[ 0 = \varphi(e_1)
   = (c_1\varphi_1+c_2\varphi_2+c_3\varphi_3+c_4\varphi_4+c_5\varphi_5)(e_1)
   = c_1. \]
Similarly, $e_2\in U$ and thus $c_2=0$. Hence
$\varphi=c_3\varphi_3+c_4\varphi_4+c_5\varphi_5$. Thus
$\varphi\in\spn(\varphi_3,\varphi_4,\varphi_5)$, which shows that
$U^0\subseteq\spn(\varphi_3,\varphi_4,\varphi_5)$.

Thus $U^0=\spn(\varphi_3,\varphi_4,\varphi_5)$.
\end{example}

\begin{result}[The annihilator is a subspace]\label{ax:3.124}
Suppose $U\subseteq V$. Then $U^0$ is a subspace of $V'$.
\end{result}

\begin{proof}
Note that $0\in U^0$ (here $0$ is the zero linear functional on $V$) because the
zero linear functional applied to every vector in $U$ equals $0\in\F$.

Suppose $\varphi,\psi\in U^0$. Thus $\varphi,\psi\in V'$ and
$\varphi(u)=\psi(u)=0$ for every $u\in U$. If $u\in U$, then
\[ (\varphi+\psi)(u) = \varphi(u)+\psi(u) = 0+0 = 0. \]
Thus $\varphi+\psi\in U^0$.

Similarly, $U^0$ is closed under scalar multiplication. Thus \ref{ax:1.34}
implies that $U^0$ is a subspace of $V'$.
\end{proof}

The next result shows that $\dim U^0$ is the difference of $\dim V$ and
$\dim U$. For example, this shows that if $U$ is a two-dimensional subspace of
$\R^5$, then $U^0$ is a three-dimensional subspace of $(\R^5)'$, as in Example
\ref{ax:3.123}.

The next result can be proved following the pattern of Example \ref{ax:3.123}:
choose a basis $u_1,\dots,u_m$ of $U$, extend to a basis
$u_1,\dots,u_m,\dots,u_n$ of $V$, let $\varphi_1,\dots,\varphi_m,\dots,\varphi_n$
be the dual basis of $V'$, and then show that $\varphi_{m+1},\dots,\varphi_n$ is
a basis of $U^0$, which implies the desired result. You should construct the
proof just outlined, even though a slicker proof is presented here.

\begin{result}[Dimension of the annihilator]\label{ax:3.125}
Suppose $V$ is finite-dimensional and $U$ is a subspace of $V$. Then
\[ \dim U^0 = \dim V-\dim U. \]
\end{result}

\begin{proof}
Let $i\in\Lin(U,V)$ be the inclusion map defined by $i(u)=u$ for each $u\in U$.
Thus $i'$ is a linear map from $V'$ to $U'$. The fundamental theorem of linear
maps (\ref{ax:3.21}) applied to $i'$ shows that
\[ \dim\range i'+\dim\nul i' = \dim V'. \]
However, $\nul i'=U^0$ (as can be seen by thinking about the definitions) and
$\dim V'=\dim V$ (by \ref{ax:3.111}), so we can rewrite the equation above
as\refstepcounter{theorem}
\begin{equation*}
\dim\range i'+\dim U^0 = \dim V. \tag*{\thetheorem}\label{ax:3.126}
\end{equation*}

If $\varphi\in U'$, then $\varphi$ can be extended to a linear functional $\psi$
on $V$ (see, for example, Exercise~13 in Section~3A). The definition of $i'$
shows that $i'(\psi)=\varphi$. Thus $\varphi\in\range i'$, which implies that
$\range i'=U'$. Hence
\[ \dim\range i' = \dim U' = \dim U, \]
and then \ref{ax:3.126} becomes the equation $\dim U+\dim U^0=\dim V$, as
desired.
\end{proof}

The next result can be a useful tool to show that a subspace is as big as
possible---see (a)---or to show that a subspace is as small as possible---see
(b).

\begin{result}[Condition for the annihilator to equal $\{0\}$ or the whole
space]\label{ax:3.127}
Suppose $V$ is finite-dimensional and $U$ is a subspace of $V$. Then
\begin{parts}
\item $U^0=\{0\}\iff U=V$;
\item $U^0=V'\iff U=\{0\}$.
\end{parts}
\end{result}

\begin{proof}
To prove (a), we have
\begin{align*}
U^0=\{0\} &\iff \dim U^0=0\\
          &\iff \dim U=\dim V\\
          &\iff U=V,
\end{align*}
where the second equivalence follows from \ref{ax:3.125} and the third
equivalence follows from \ref{ax:2.39}.

Similarly, to prove (b) we have
\begin{align*}
U^0=V' &\iff \dim U^0=\dim V'\\
       &\iff \dim U^0=\dim V\\
       &\iff \dim U=0\\
       &\iff U=\{0\},
\end{align*}
where one direction of the first equivalence follows from \ref{ax:2.39}, the
second equivalence follows from \ref{ax:3.111}, and the third equivalence
follows from \ref{ax:3.125}.
\end{proof}

The proof of (a) in the next result does not use the hypothesis that $V$ and
$W$ are finite-dimensional.

\begin{result}[The null space of $T'$]\label{ax:3.128}
Suppose $V$ and $W$ are finite-dimensional and $T\in\Lin(V,W)$. Then
\begin{parts}
\item $\nul T'=(\range T)^0$;
\item $\dim\nul T'=\dim\nul T+\dim W-\dim V$.
\end{parts}
\end{result}

\begin{proof}\leavevmode
\begin{parts}
\item First suppose $\varphi\in\nul T'$. Thus $0=T'(\varphi)=\varphi\circ T$.
  Hence
  \[ 0 = (\varphi\circ T)(v) = \varphi(Tv) \quad\text{for every } v\in V. \]
  Thus $\varphi\in(\range T)^0$. This implies that
  $\nul T'\subseteq(\range T)^0$.

  To prove the inclusion in the opposite direction, now suppose
  $\varphi\in(\range T)^0$. Thus $\varphi(Tv)=0$ for every vector $v\in V$.
  Hence $0=\varphi\circ T=T'(\varphi)$. In other words, $\varphi\in\nul T'$,
  which shows that $(\range T)^0\subseteq\nul T'$, completing the proof of (a).
\item We have
  \begin{align*}
  \dim\nul T' &= \dim(\range T)^0\\
              &= \dim W-\dim\range T\\
              &= \dim W-(\dim V-\dim\nul T)\\
              &= \dim\nul T+\dim W-\dim V,
  \end{align*}
  where the first equality comes from (a), the second equality comes from
  \ref{ax:3.125}, and the third equality comes from the fundamental theorem of
  linear maps (\ref{ax:3.21}). \qedhere
\end{parts}
\end{proof}

The next result can be useful because sometimes it is easier to verify that $T'$
is injective than to show directly that $T$ is surjective.

\begin{result}[$T$ surjective is equivalent to $T'$ injective]\label{ax:3.129}
Suppose $V$ and $W$ are finite-dimensional and $T\in\Lin(V,W)$. Then
\[ T \text{ is surjective} \iff T' \text{ is injective}. \]
\end{result}

\begin{proof}
We have
\begin{align*}
T\in\Lin(V,W)\text{ is surjective} &\iff \range T=W\\
  &\iff (\range T)^0=\{0\}\\
  &\iff \nul T'=\{0\}\\
  &\iff T'\text{ is injective},
\end{align*}
where the second equivalence comes from \ref{ax:3.127}(a) and the third
equivalence comes from \ref{ax:3.128}(a).
\end{proof}

\begin{result}[The range of $T'$]\label{ax:3.130}
Suppose $V$ and $W$ are finite-dimensional and $T\in\Lin(V,W)$. Then
\begin{parts}
\item $\dim\range T'=\dim\range T$;
\item $\range T'=(\nul T)^0$.
\end{parts}
\end{result}

\begin{proof}\leavevmode
\begin{parts}
\item We have
  \begin{align*}
  \dim\range T' &= \dim W'-\dim\nul T'\\
                &= \dim W-\dim(\range T)^0\\
                &= \dim\range T,
  \end{align*}
  where the first equality comes from \ref{ax:3.21}, the second equality comes
  from \ref{ax:3.111} and \ref{ax:3.128}(a), and the third equality comes from
  \ref{ax:3.125}.
\item First suppose $\varphi\in\range T'$. Thus there exists $\psi\in W'$ such
  that $\varphi=T'(\psi)$. If $v\in\nul T$, then
  \[ \varphi(v) = \bigl(T'(\psi)\bigr)v = (\psi\circ T)(v) = \psi(Tv)
     = \psi(0) = 0. \]
  Hence $\varphi\in(\nul T)^0$. This implies that $\range T'\subseteq(\nul T)^0$.

  We will complete the proof by showing that $\range T'$ and $(\nul T)^0$ have
  the same dimension. To do this, note that
  \begin{align*}
  \dim\range T' &= \dim\range T\\
                &= \dim V-\dim\nul T\\
                &= \dim(\nul T)^0,
  \end{align*}
  where the first equality comes from (a), the second equality comes from
  \ref{ax:3.21}, and the third equality comes from \ref{ax:3.125}. \qedhere
\end{parts}
\end{proof}

The next result should be compared to \ref{ax:3.129}.

\begin{result}[$T$ injective is equivalent to $T'$ surjective]\label{ax:3.131}
Suppose $V$ and $W$ are finite-dimensional and $T\in\Lin(V,W)$. Then
\[ T \text{ is injective} \iff T' \text{ is surjective}. \]
\end{result}

\begin{proof}
We have
\begin{align*}
T\text{ is injective} &\iff \nul T=\{0\}\\
  &\iff (\nul T)^0=V'\\
  &\iff \range T'=V',
\end{align*}
where the second equivalence follows from \ref{ax:3.127}(b) and the third
equivalence follows from \ref{ax:3.130}(b).
\end{proof}

\subsection{Matrix of Dual of Linear Map}

The setting for the next result is the assumption that we have a basis
$v_1,\dots,v_n$ of $V$, along with its dual basis $\varphi_1,\dots,\varphi_n$ of
$V'$. We also have a basis $w_1,\dots,w_m$ of $W$, along with its dual basis
$\psi_1,\dots,\psi_m$ of $W'$. Thus $\Mat(T)$ is computed with respect to the
bases just mentioned of $V$ and $W$, and $\Mat(T')$ is computed with respect to
the dual bases just mentioned of $W'$ and $V'$. Using these bases gives the
following pretty result.

\begin{result}[Matrix of $T'$ is transpose of matrix of $T$]\label{ax:3.132}
Suppose $V$ and $W$ are finite-dimensional and $T\in\Lin(V,W)$. Then
\[ \Mat(T') = \bigl(\Mat(T)\bigr)^{\mathrm{t}}. \]
\end{result}

\begin{proof}
Let $A=\Mat(T)$ and $C=\Mat(T')$. Suppose $1\le j\le m$ and $1\le k\le n$.

From the definition of $\Mat(T')$ we have
\[ T'(\psi_j) = \sum_{r=1}^n C_{r,j}\varphi_r. \]
The left side of the equation above equals $\psi_j\circ T$. Thus applying both
sides of the equation above to $v_k$ gives
\begin{align*}
(\psi_j\circ T)(v_k) &= \sum_{r=1}^n C_{r,j}\varphi_r(v_k)\\
                     &= C_{k,j}.
\end{align*}
We also have
\begin{align*}
(\psi_j\circ T)(v_k) &= \psi_j(Tv_k)\\
  &= \psi_j\Bigl(\sum_{r=1}^m A_{r,k}w_r\Bigr)\\
  &= \sum_{r=1}^m A_{r,k}\psi_j(w_r)\\
  &= A_{j,k}.
\end{align*}
Comparing the last line of the last two sets of equations, we have
$C_{k,j}=A_{j,k}$. Thus $C=A^{\mathrm{t}}$. In other words,
$\Mat(T')=\bigl(\Mat(T)\bigr)^{\mathrm{t}}$, as desired.
\end{proof}

Now we use duality to give an alternative proof that the column rank of a matrix
equals the row rank of the matrix. This result was previously proved using
different tools---see \ref{ax:3.57}.

\begin{result}[Column rank equals row rank]\label{ax:3.133}
Suppose $A\in\F^{m,n}$. Then the column rank of $A$ equals the row rank of $A$.
\end{result}

\begin{proof}
Define $T\colon\F^{n,1}\to\F^{m,1}$ by $Tx=Ax$. Thus $\Mat(T)=A$, where
$\Mat(T)$ is computed with respect to the standard bases of $\F^{n,1}$ and
$\F^{m,1}$. Now
\begin{align*}
\text{column rank of } A &= \text{column rank of } \Mat(T)\\
  &= \dim\range T\\
  &= \dim\range T'\\
  &= \text{column rank of } \Mat(T')\\
  &= \text{column rank of } A^{\mathrm{t}}\\
  &= \text{row rank of } A,
\end{align*}
where the second equality comes from \ref{ax:3.78}, the third equality comes
from \ref{ax:3.130}(a), the fourth equality comes from \ref{ax:3.78}, the fifth
equality comes from \ref{ax:3.132}, and the last equality follows from the
definitions of row and column rank.
\end{proof}

See Exercise~8 in Section~7A for another alternative proof of the result above.

% ------------------------------------------------------------------ exercises
\exercisesheading{3F}

\begin{exercise}
Explain why each linear functional is surjective or is the zero map.
\end{exercise}
\begin{solution}
\solnote{Manual Remark}{For the rest of the solutions manual, unless otherwise stated, we
will be using the Kronecker delta notation, i.e., we define $\delta_{ij}=1$ if
$i=j$ and $\delta_{ij}=0$ if $i\ne j$.}

Let $\varphi\in V'$. If $\varphi=0$, then it is the zero map. Otherwise, there
exists $v\in V$ such that $\varphi(v)\ne0$, say $\varphi(v)=c\ne0$. Then for any
$x\in\F$, we have
\[ \varphi\Bigl(\frac{x}{c}v\Bigr) = \frac{x}{c}\varphi(v) = \frac{x}{c}c = x. \]
Hence, $\varphi$ is surjective.
\end{solution}

\begin{exercise}
Give three distinct examples of linear functionals on $\R^{[0,1]}$.
\end{exercise}
\begin{solution}
Consider the linear functionals on $\R^{[0,1]}$ defined by
\[ \varphi_1(f)=f(0), \quad \varphi_2(f)=f(1), \quad
   \varphi_3(f)=f\Bigl(1/2\Bigr). \qedhere \]
\end{solution}

\begin{exercise}
Suppose $V$ is finite-dimensional and $v\in V$ with $v\ne0$. Prove that there
exists $\varphi\in V'$ such that $\varphi(v)=1$.
\end{exercise}
\begin{solution}
Since $V$ is finite-dimensional, we can extend $v$ to a basis of $V$, say
$v,v_2,\dots,v_n$. Define $\varphi\in V'$ by setting $\varphi(v)=1$ and
$\varphi(v_i)=0$ for $i=2,\dots,n$. This defines a linear functional because it
is defined on a basis and extended linearly. Clearly, $\varphi(v)=1$.
\end{solution}

\begin{exercise}
Suppose $V$ is finite-dimensional and $U$ is a subspace of $V$ such that
$U\ne V$. Prove that there exists $\varphi\in V'$ such that $\varphi(u)=0$ for
every $u\in U$ but $\varphi\ne0$.
\end{exercise}
\begin{solution}
Since $U\ne V$, there exists $v_1\in V$ such that $v_1\notin U$. Extend a basis
of $U$ to a basis of $V$, say $u_1,\dots,u_k,v_1,\dots,v_m$. Define
$\varphi\in V'$ by setting $\varphi(u_i)=0$ for all $i=1,\dots,k$ and
$\varphi(v_j)=0$ for all $j=2,\dots,m$ but $\varphi(v_1)=1$. This defines a
linear functional because it is defined on a basis and extended linearly.
Clearly, $\varphi(u)=0$ for all $u\in U$ and $\varphi\ne0$.
\end{solution}

\begin{exercise}
Suppose $T\in\Lin(V,W)$ and $w_1,\dots,w_m$ is a basis of $\range T$. Hence for
each $v\in V$, there exist unique numbers $\varphi_1(v),\dots,\varphi_m(v)$ such
that
\[ Tv = \varphi_1(v)w_1+\cdots+\varphi_m(v)w_m, \]
thus defining functions $\varphi_1,\dots,\varphi_m$ from $V$ to $\F$. Show that
each of the functions $\varphi_1,\dots,\varphi_m$ is a linear functional on $V$.
\end{exercise}
\begin{solution}
Let $v\in V$. Since $w_1,\dots,w_m$ is a basis of $\range T$, there exist unique
scalars $\varphi_1(v),\dots,\varphi_m(v)$ such that
\[ Tv = \varphi_1(v)w_1+\cdots+\varphi_m(v)w_m. \]
To show that each $\varphi_i$ is a linear functional, let $v,u\in V$ and
$a,b\in\F$. Then
\begin{multline*}
T(av+bu) = aTv+bTu\\
 = a\bigl(\varphi_1(v)w_1+\cdots+\varphi_m(v)w_m\bigr)
   +b\bigl(\varphi_1(u)w_1+\cdots+\varphi_m(u)w_m\bigr).
\end{multline*}
By the uniqueness of the representation in terms of the basis $w_1,\dots,w_m$,
we must have
\[ \varphi_i(av+bu) = a\varphi_i(v)+b\varphi_i(u) \quad\text{for each }
   i=1,\dots,m. \]
Hence, each $\varphi_i$ is a linear functional on $V$.
\end{solution}

\begin{exercise}
Suppose $\varphi,\beta\in V'$. Prove that $\nul\varphi\subseteq\nul\beta$ if and
only if there exists $c\in\F$ such that $\beta=c\varphi$.
\end{exercise}
\begin{solution}
Suppose $\nul\varphi\subseteq\nul\beta$. If $\varphi=0$, then $\nul\varphi=V$,
and hence $\nul\beta=V$, which implies $\beta=0$. In this case, we can take
$c=0$ and have $\beta=c\varphi$. Now assume $\varphi\ne0$. Pick $v_0\in V$ such
that $\varphi(v_0)\ne0$. Define
\[ c = \frac{\beta(v_0)}{\varphi(v_0)}. \]
For any $v\in V$, consider
\[ v-\frac{\varphi(v)}{\varphi(v_0)}v_0. \]
This vector is in $\nul\varphi$, so it is also in $\nul\beta$. Hence,
\[ \beta\Bigl(v-\frac{\varphi(v)}{\varphi(v_0)}v_0\Bigr) = 0
   \implies \beta(v) = \frac{\varphi(v)}{\varphi(v_0)}\beta(v_0) = c\varphi(v). \]
Therefore, $\beta=c\varphi$.

Conversely, if $\beta=c\varphi$, then clearly $\nul\varphi\subseteq\nul\beta$
because if $\varphi(v)=0$, then $\beta(v)=c\varphi(v)=0$.
\end{solution}

\begin{exercise}
Suppose that $V_1,\dots,V_m$ are vector spaces. Prove that
$(V_1\times\cdots\times V_m)'$ and ${V_1}'\times\cdots\times{V_m}'$ are
isomorphic vector spaces.
\end{exercise}
\begin{solution}
Define a map
$\Phi\colon(V_1\times\cdots\times V_m)'\to{V_1}'\times\cdots\times{V_m}'$ by
\[ \Phi(\varphi) = (\varphi_1,\dots,\varphi_m), \]
where $\varphi_i\in{V_i}'$ is defined by
$\varphi_i(v_i)=\varphi(0,\dots,0,v_i,0,\dots,0)$, with $v_i$ in the
$i$th position.

First, we show that $\Phi$ is linear. Let
$\varphi,\psi\in(V_1\times\cdots\times V_m)'$ and $a,b\in\F$. Then
\begin{align*}
\Phi(a\varphi+b\psi)
 &= \bigl((a\varphi+b\psi)_1,\dots,(a\varphi+b\psi)_m\bigr)\\
 &= (a\varphi_1+b\psi_1,\dots,a\varphi_m+b\psi_m) = a\Phi(\varphi)+b\Phi(\psi).
\end{align*}

Next, we show that $\Phi$ is bijective. For injectivity, suppose
$\Phi(\varphi)=0$. Then $\varphi_i=0$ for all $i$, which implies
$\varphi(0,\dots,0,v_i,0,\dots,0)=0$ for all $v_i\in V_i$. By linearity,
$\varphi$ must be zero on all of $V_1\times\cdots\times V_m$, so $\varphi=0$.
For surjectivity, given
$(\varphi_1,\dots,\varphi_m)\in{V_1}'\times\cdots\times{V_m}'$, define
$\varphi\in(V_1\times\cdots\times V_m)'$ by
\[ \varphi(v_1,\dots,v_m) = \varphi_1(v_1)+\cdots+\varphi_m(v_m). \]
Then $\Phi(\varphi)=(\varphi_1,\dots,\varphi_m)$, so $\Phi$ is surjective.

Therefore, $\Phi$ is a linear isomorphism, and $(V_1\times\cdots\times V_m)'$
and ${V_1}'\times\cdots\times{V_m}'$ are isomorphic vector spaces.
\end{solution}

\begin{exercise}
Suppose $v_1,\dots,v_n$ is a basis of $V$ and $\varphi_1,\dots,\varphi_n$ is the
dual basis of $V'$. Define $\Gamma\colon V\to\F^n$ and $\Lambda\colon\F^n\to V$
by
\[ \Gamma(v) = \bigl(\varphi_1(v),\dots,\varphi_n(v)\bigr) \quad\text{and}\quad
   \Lambda(a_1,\dots,a_n) = a_1v_1+\cdots+a_nv_n. \]
Explain why $\Gamma$ and $\Lambda$ are inverses of each other.
\end{exercise}
\begin{solution}
To show that $\Gamma$ and $\Lambda$ are inverses of each other, we need to check
that $\Gamma\bigl(\Lambda(a_1,\dots,a_n)\bigr)=(a_1,\dots,a_n)$ for all
$(a_1,\dots,a_n)\in\F^n$ and $\Lambda\bigl(\Gamma(v)\bigr)=v$ for all $v\in V$.

First, let $(a_1,\dots,a_n)\in\F^n$. Then
\begin{align*}
\Gamma\bigl(\Lambda(a_1,\dots,a_n)\bigr) &= \Gamma(a_1v_1+\cdots+a_nv_n)\\
 &= \bigl(\varphi_1(a_1v_1+\cdots+a_nv_n),\dots,
    \varphi_n(a_1v_1+\cdots+a_nv_n)\bigr).
\end{align*}
By the definition of the dual basis, we have $\varphi_i(v_j)=\delta_{ij}$ for all
$i,j$. Therefore,
\[ \varphi_i(a_1v_1+\cdots+a_nv_n) = a_i. \]
Therefore,
\[ \Gamma\bigl(\Lambda(a_1,\dots,a_n)\bigr) = (a_1,\dots,a_n). \]
Next, let $v\in V$. Write $v=b_1v_1+\cdots+b_nv_n$ for some scalars
$b_1,\dots,b_n$. Then
\[ \Lambda\bigl(\Gamma(v)\bigr) = \Lambda\bigl(\varphi_1(v),\dots,\varphi_n(v)\bigr)
   = \Lambda(b_1,\dots,b_n) = b_1v_1+\cdots+b_nv_n = v. \]

Hence, $\Gamma$ and $\Lambda$ are inverses of each other.
\end{solution}

\begin{exercise}
Suppose $m$ is a positive integer. Show that the dual basis of the basis
$1,x,\dots,x^m$ of $\Poly_m(\R)$ is $\varphi_0,\varphi_1,\dots,\varphi_m$, where
\[ \varphi_k(p) = \frac{p^{(k)}(0)}{k!}. \]
\exnote{Here $p^{(k)}$ denotes the $k^{\text{th}}$ derivative of $p$, with the
understanding that the $0^{\text{th}}$ derivative of $p$ is $p$.}
\end{exercise}
\begin{solution}
Let $p\in\Poly_m(\R)$. We can write $p(x)=a_0+a_1x+\cdots+a_mx^m$ for some
scalars $a_0,\dots,a_m$. Then for each $k=0,1,\dots,m$, we have
\begin{align*}
p(x) &= a_0+a_1x+a_2x^2+\cdots+a_mx^m\\
p'(x) &= a_1+2a_2x+\cdots+ma_mx^{m-1}\\
p''(x) &= 2a_2+\cdots+m(m-1)a_mx^{m-2}\\
&\;\;\vdots\\
p^{(k)}(x) &= k!a_k+\cdots+\frac{m!}{(m-k)!}a_mx^{m-k}
\end{align*}
Then
\[ \varphi_k(p) = \frac{p^{(k)}(0)}{k!} = a_k. \]
Therefore, $\varphi_k(x^j)=\delta_{kj}$ for all $0\le k,j\le m$, which shows
that $\varphi_0,\varphi_1,\dots,\varphi_m$ is indeed the dual basis of
$1,x,\dots,x^m$.
\end{solution}

\begin{exercise}
Suppose $m$ is a positive integer.
\begin{parts}
\item Show that $1,x-5,\dots,(x-5)^m$ is a basis of $\Poly_m(\R)$.
\item What is the dual basis of the basis in (a)?
\end{parts}
\end{exercise}
\begin{solution}\leavevmode
\begin{parts}
\item Note that $1,x-5,\dots,(x-5)^m$ is a list of $m+1$ polynomials in
  $\Poly_m(\R)$ with differing degrees, and since $\Poly_m(\R)$ has dimension
  $m+1$, this list forms a basis of $\Poly_m(\R)$.
\item Substitute $y=x-5$. Then the basis $1,x-5,\dots,(x-5)^m$ becomes
  $1,y,\dots,y^m$. By Exercise~3F9, the dual basis of $1,y,\dots,y^m$ is
  $\varphi_0,\varphi_1,\dots,\varphi_m$, where
  \[ \varphi_k(p) = \frac{p^{(k)}(0)}{k!}. \]
  Since $y=x-5$, then
  \[ \varphi_k(p) = \frac{p^{(k)}(5)}{k!}. \qedhere \]
\end{parts}
\end{solution}

\begin{exercise}
Suppose $v_1,\dots,v_n$ is a basis of $V$ and $\varphi_1,\dots,\varphi_n$ is the
corresponding dual basis of $V'$. Suppose $\psi\in V'$. Prove that
\[ \psi = \psi(v_1)\varphi_1+\cdots+\psi(v_n)\varphi_n. \]
\end{exercise}
\begin{solution}
Let $v\in V$. Since $\varphi_1,\dots,\varphi_n$ is the dual basis of
$v_1,\dots,v_n$, we have $\varphi_i(v_j)=\delta_{ij}$ for all $1\le i,j\le n$.
Let $v=\alpha_1v_1+\cdots+\alpha_nv_n$ for some scalars
$\alpha_1,\dots,\alpha_n$. Then
\begin{align*}
\bigl(\psi(v_1)\varphi_1+\cdots+\psi(v_n)\varphi_n\bigr)(v)
 &= \psi(v_1)\varphi_1(v)+\cdots+\psi(v_n)\varphi_n(v)\\
 &= \psi(v_1)\alpha_1+\cdots+\psi(v_n)\alpha_n\\
 &= \psi(\alpha_1v_1+\cdots+\alpha_nv_n)\\
 &= \psi(v)
\end{align*}
Since this holds for all $v\in V$, we conclude that
\[ \psi = \psi(v_1)\varphi_1+\cdots+\psi(v_n)\varphi_n. \]

\emph{Alternate proof.} \emph{Civil Service Pigeon.} Because
$\varphi_1,\dots,\varphi_n$ is the dual basis, $\varphi_i(v_j)=\delta_{i,j}$.
For each basis vector $v_j$,
\[ \Bigl(\sum_{i=1}^n\psi(v_i)\varphi_i\Bigr)(v_j)
   = \sum_{i=1}^n\psi(v_i)\varphi_i(v_j) = \psi(v_j). \]
Since the two linear functionals agree on a basis of $V$, they are equal on all
of $V$.
\end{solution}

\begin{exercise}
Suppose $S,T\in\Lin(V,W)$.
\begin{parts}
\item Prove that $(S+T)'=S'+T'$.
\item Prove that $(\lambda T)'=\lambda T'$ for all $\lambda\in\F$.
\end{parts}
\exnote{This exercise asks you to verify \textup{(a)} and \textup{(b)} in
\ref{ax:3.120}.}
\end{exercise}
\begin{solution}\leavevmode
\begin{parts}
\item Let $\psi\in W'$. Then for any $v\in V$, we have
  \begin{align*}
  \bigl((S+T)'(\psi)\bigr)(v) &= \psi\bigl((S+T)(v)\bigr)\\
   &= \psi\bigl(S(v)+T(v)\bigr)\\
   &= \psi\bigl(S(v)\bigr)+\psi\bigl(T(v)\bigr)\\
   &= \bigl(S'(\psi)\bigr)(v)+\bigl(T'(\psi)\bigr)(v)\\
   &= \bigl(S'(\psi)+T'(\psi)\bigr)(v).
  \end{align*}
  Since this holds for all $v\in V$, we conclude that
  \[ (S+T)'(\psi) = S'(\psi)+T'(\psi). \]
  Since this holds for all $\psi\in W'$, we have
  \[ (S+T)' = S'+T'. \]
\item Let $\psi\in W'$. Then for any $v\in V$, we have
  \begin{align*}
  \bigl((\lambda T)'(\psi)\bigr)(v) &= \psi\bigl((\lambda T)(v)\bigr)\\
   &= \psi\bigl(\lambda T(v)\bigr)\\
   &= \lambda\psi\bigl(T(v)\bigr)\\
   &= \bigl(\lambda T'(\psi)\bigr)(v).
  \end{align*}
  Since this holds for all $v\in V$, we conclude that
  \[ (\lambda T)'(\psi) = \lambda T'(\psi). \]
  Since this holds for all $\psi\in W'$, we have
  \[ (\lambda T)' = \lambda T'. \qedhere \]
\end{parts}
\end{solution}

\begin{exercise}
Show that the dual map of the identity operator on $V$ is the identity operator
on $V'$.
\end{exercise}
\begin{solution}
Let $\psi\in V'$. Let $I'$ denote the dual map of the identity operator $I$ on
$V$. Then for any $v\in V$, we have
\begin{align*}
\bigl(I'(\psi)\bigr)(v) &= \psi\bigl(I(v)\bigr)\\
 &= \psi(v).
\end{align*}
Since this holds for all $v\in V$, we conclude that
\[ I'(\psi) = \psi. \]
Since this holds for all $\psi\in V'$, then $I'$ must be the identity operator
on $V'$.
\end{solution}

\begin{exercise}
Define $T\colon\R^3\to\R^2$ by
\[ T(x,y,z) = (4x+5y+6z,7x+8y+9z). \]
Suppose $\varphi_1,\varphi_2$ denotes the dual basis of the standard basis of
$\R^2$ and $\psi_1,\psi_2,\psi_3$ denotes the dual basis of the standard basis
of $\R^3$.
\begin{parts}
\item Describe the linear functionals $T'(\varphi_1)$ and $T'(\varphi_2)$.
\item Write $T'(\varphi_1)$ and $T'(\varphi_2)$ as linear combinations of
  $\psi_1,\psi_2,\psi_3$.
\end{parts}
\end{exercise}
\begin{solution}\leavevmode
\begin{parts}
\item Let $(x,y,z)\in\R^3$. Then we have
  \begin{align*}
  \bigl(T'(\varphi_1)\bigr)\bigl((x,y,z)\bigr)
   &= \varphi_1\bigl(T(x,y,z)\bigr)\\
   &= \varphi_1(4x+5y+6z,7x+8y+9z)\\
   &= 4x+5y+6z,\\
  \bigl(T'(\varphi_2)\bigr)\bigl((x,y,z)\bigr)
   &= \varphi_2\bigl(T(x,y,z)\bigr)\\
   &= \varphi_2(4x+5y+6z,7x+8y+9z)\\
   &= 7x+8y+9z.
  \end{align*}
\item We can write $T'(\varphi_1)$ and $T'(\varphi_2)$ as linear combinations
  of $\psi_1,\psi_2,\psi_3$ as follows:
  \begin{align*}
  T'(\varphi_1) &= 4\psi_1+5\psi_2+6\psi_3,\\
  T'(\varphi_2) &= 7\psi_1+8\psi_2+9\psi_3. \qedhere
  \end{align*}
\end{parts}
\end{solution}

\begin{exercise}
Define $T\colon\Poly(\R)\to\Poly(\R)$ by
\[ (Tp)(x) = x^2p(x)+p''(x) \]
for each $x\in\R$.
\begin{parts}
\item Suppose $\varphi\in\Poly(\R)'$ is defined by $\varphi(p)=p'(4)$. Describe
  the linear functional $T'(\varphi)$ on $\Poly(\R)$.
\item Suppose $\varphi\in\Poly(\R)'$ is defined by $\varphi(p)=\int_0^1p$.
  Evaluate $\bigl(T'(\varphi)\bigr)(x^3)$.
\end{parts}
\end{exercise}
\begin{solution}\leavevmode
\begin{parts}
\item Let $p\in\Poly(\R)$. Then we have
  \begin{align*}
  \bigl(T'(\varphi)\bigr)(p) &= \varphi\bigl(T(p)\bigr)\\
   &= (Tp)'(4)\\
   &= \bigl(2xp(x)+x^2p'(x)+p'''(x)\bigr)\big|_{x=4}\\
   &= 8p(4)+16p'(4)+p'''(4).
  \end{align*}
\item Let $p(x)=x^3$. Then we have
  \begin{align*}
  \bigl(T'(\varphi)\bigr)(x^3) &= \varphi\bigl(T(x^3)\bigr)\\
   &= \int_0^1 T(x^3)\,\mathrm{d}x\\
   &= \int_0^1 \bigl(x^2\cdot x^3+(x^3)''\bigr)\,\mathrm{d}x\\
   &= \int_0^1 (x^5+6x)\,\mathrm{d}x\\
   &= 1/6+18/6\\
   &= 19/6. \qedhere
  \end{align*}
\end{parts}
\end{solution}

\begin{exercise}
Suppose $W$ is finite-dimensional and $T\in\Lin(V,W)$. Prove that
\[ T'=0 \iff T=0. \]
\end{exercise}
\begin{solution}
Suppose $T'=0$. Then for every $\varphi\in W'$ and every $v\in V$, we have
\[ \bigl(T'(\varphi)\bigr)(v) = \varphi\bigl(T(v)\bigr) = 0. \]
If $T(v)\ne0$ for some $v\in V$, then $T(v)\in W$. Since $W$ is
finite-dimensional, then by Exercise~3F3, there exists $\psi\in W'$ such that
$\psi\bigl(T(v)\bigr)\ne0$, which contradicts the fact that
$\bigl(T'(\psi)\bigr)(v)=\psi\bigl(T(v)\bigr)=0$. Hence, $T(v)=0$ for every
$v\in V$, and so $T=0$.

Conversely, suppose $T=0$. Then for every $\varphi\in W'$ and every $v\in V$, we
have
\[ \bigl(T'(\varphi)\bigr)(v) = \varphi\bigl(T(v)\bigr) = \varphi(0) = 0. \]
Hence, $T'=0$.
\end{solution}

\begin{exercise}
Suppose $V$ and $W$ are finite-dimensional and $T\in\Lin(V,W)$. Prove that $T$
is invertible if and only if $T'\in\Lin(W',V')$ is invertible.
\end{exercise}
\begin{solution}
Suppose $T$ is invertible. Then there exists $S\in\Lin(W,V)$ such that
$ST=I_V$ and $TS=I_W$. Taking the dual of both sides, we get $T'S'=I_{W'}$ and
$S'T'=I_{V'}$, which shows that $T'$ is invertible with inverse $S'$.

Conversely, suppose $T'$ is invertible. Then $T'$ is an isomorphism. Then
\[ \dim V = \dim V' = \dim W' = \dim W. \]
By injectivity of $T'$, we have $\nul T'=\{0\}$. Since
$\nul T'=(\range T)^0$, we have $(\range T)^0=\{0\}$ so that $\range T=W$.
Hence, $T$ is surjective. Since $\dim V=\dim W$, we conclude that $T$ is
injective as well. Therefore, $T$ is invertible.
\end{solution}

\begin{exercise}
Suppose $V$ and $W$ are finite-dimensional. Prove that the map that takes
$T\in\Lin(V,W)$ to $T'\in\Lin(W',V')$ is an isomorphism of $\Lin(V,W)$ onto
$\Lin(W',V')$.
\end{exercise}
\begin{solution}
Let $\Phi\colon\Lin(V,W)\to\Lin(W',V')$ be the map defined by $\Phi(T)=T'$. We
need to show that $\Phi$ is linear, injective, and surjective.

Linearity: For $T_1,T_2\in\Lin(V,W)$ and $a\in\F$, we have
\[ \Phi(T_1+T_2) = (T_1+T_2)' = {T_1}'+{T_2}' = \Phi(T_1)+\Phi(T_2), \]
and
\[ \Phi(aT_1) = (aT_1)' = a{T_1}' = a\Phi(T_1). \]
Hence, $\Phi$ is linear.

To see that $\Phi$ is injective, suppose $\Phi(T)=T'=0$. By Exercise~3F16, this
implies $T=0$. Hence, $\Phi$ is injective.

To see that $\Phi$ is surjective, note that since $V$ and $W$ are
finite-dimensional,
$\dim\Lin(V,W)=\dim V\cdot\dim W=\dim W'\cdot\dim V'=\dim\Lin(W',V')$. A linear
map between finite-dimensional vector spaces of the same dimension that is
injective is also surjective. Hence, $\Phi$ is surjective and therefore an
isomorphism.
\end{solution}

\begin{exercise}
Suppose $U\subseteq V$. Explain why
\[ U^0 = \{\varphi\in V' : U\subseteq\nul\varphi\}. \]
\end{exercise}
\begin{solution}
By definition, the annihilator $U^0$ of $U$ is the set of all linear
functionals $\varphi\in V'$ such that $\varphi(u)=0$ for all $u\in U$.
Equivalently, this means that $U\subseteq\nul\varphi$.
\end{solution}

\begin{exercise}
Suppose $V$ is finite-dimensional and $U$ is a subspace of $V$. Show that
\[ U = \{v\in V : \varphi(v)=0 \text{ for every } \varphi\in U^0\}. \]
\end{exercise}
\begin{solution}
Let $S=\{v\in V \mid \varphi(v)=0 \text{ for every } \varphi\in U^0\}$. We need
to show that $S=U$.

First, if $v\in U$, then for any $\varphi\in U^0$, we have $\varphi(v)=0$ by
definition of the annihilator. Hence, $v\in S$, so $U\subseteq S$.

Conversely, suppose $v\in S$. If $v\notin U$, then by Exercise~3F4, there
exists a linear functional $\varphi\in V'$ such that $\varphi(u)=0$ for all
$u\in U$ and $\varphi(v)\ne0$. But this $\varphi$ is in $U^0$, which
contradicts $v\in S$. Hence, $v\in U$, and therefore $S\subseteq U$.

Combining both inclusions, we get $S=U$, as desired.
\end{solution}

\begin{exercise}
Suppose $V$ is finite-dimensional and $U$ and $W$ are subspaces of $V$.
\begin{parts}
\item Prove that $W^0\subseteq U^0$ if and only if $U\subseteq W$.
\item Prove that $W^0=U^0$ if and only if $U=W$.
\end{parts}
\end{exercise}
\begin{solution}\leavevmode
\begin{parts}
\item Suppose $W^0\subseteq U^0$. If $u\in U$, then for any $\varphi\in W^0$,
  we have $\varphi(u)=0$ because $\varphi\in U^0$. Hence, by Exercise~3F20, we
  have $u\in W$, so $U\subseteq W$. Conversely, if $U\subseteq W$ and
  $\varphi\in W^0$, then $\varphi(w)=0$ for all $w\in W$, and in particular for
  all $u\in U$. Hence, $\varphi\in U^0$, so $W^0\subseteq U^0$.
\item Suppose $W^0=U^0$. Then $W^0\subseteq U^0$ and $U^0\subseteq W^0$. By
  part (a), this implies that $U\subseteq W$ and $W\subseteq U$. Hence, $U=W$.
  Conversely, if $U=W$, then clearly $U^0=W^0$.
\end{parts}
\end{solution}

\begin{exercise}
Suppose $V$ is finite-dimensional and $U$ and $W$ are subspaces of $V$.
\begin{parts}
\item Show that $(U+W)^0=U^0\cap W^0$.
\item Show that $(U\cap W)^0=U^0+W^0$.
\end{parts}
\end{exercise}
\begin{solution}\leavevmode
\begin{parts}
\item Let $\varphi\in(U+W)^0$. Then for any $u\in U$ and $w\in W$, we have
  $\varphi(u+w)=0$. In particular, taking $w=0$ gives $\varphi(u)=0$ for all
  $u\in U$, so $\varphi\in U^0$. Similarly, taking $u=0$ gives $\varphi(w)=0$
  for all $w\in W$, so $\varphi\in W^0$. Hence, $\varphi\in U^0\cap W^0$, and we
  have shown that $(U+W)^0\subseteq U^0\cap W^0$.

  Conversely, let $\varphi\in U^0\cap W^0$. Then for any $u\in U$ and $w\in W$,
  we have $\varphi(u)=0$ and $\varphi(w)=0$, so
  $\varphi(u+w)=\varphi(u)+\varphi(w)=0$. Hence, $\varphi\in(U+W)^0$, and we
  have shown that $U^0\cap W^0\subseteq(U+W)^0$. Combining both inclusions, we
  get $(U+W)^0=U^0\cap W^0$.
\item From part (a), we know that $(U+W)^0=U^0\cap W^0$. Then
  \[ \dim(U+W)^0 = \dim V-\dim(U+W). \]
  Applying the dimension formula for a sum of subspaces, we have
  \[ \dim(U^0+W^0) = \dim U^0+\dim W^0-\dim(U^0\cap W^0). \]
  Since $(U+W)^0=U^0\cap W^0$, we have
  $\dim(U^0\cap W^0)=\dim(U+W)^0=\dim V-\dim(U+W)$. Therefore,
  \[ \dim(U^0+W^0) = \dim U^0+\dim W^0-\bigl(\dim V-\dim(U+W)\bigr). \]
  Using the fact that $\dim U^0=\dim V-\dim U$ and $\dim W^0=\dim V-\dim W$, we
  get
  \begin{align*}
  \dim(U^0+W^0)
   &= \dim U^0+\dim W^0-\dim(U^0\cap W^0)\\
   &= (\dim V-\dim U)+(\dim V-\dim W)-\bigl(\dim V-\dim(U+W)\bigr)\\
   &= \dim V-\dim U+\dim V-\dim W-\dim V+\dim(U+W)\\
   &= \dim V-\bigl(\dim U+\dim W-\dim(U+W)\bigr)\\
   &= \dim V-\dim(U\cap W).
  \end{align*}
  Hence, $\dim(U^0+W^0)=\dim(U\cap W)^0$. To see equality, we just need to
  prove one inclusion. If $\varphi=\varphi_U+\varphi_W\in U^0+W^0$, then for any
  $v\in U\cap W$, we have $\varphi(v)=\varphi_U(v)+\varphi_W(v)=0+0=0$, so
  $\varphi\in(U\cap W)^0$. Therefore, $U^0+W^0\subseteq(U\cap W)^0$. Since the
  dimensions are equal, we conclude that $(U\cap W)^0=U^0+W^0$.
\end{parts}
\end{solution}

\begin{exercise}
Suppose $V$ is finite-dimensional and $\varphi_1,\dots,\varphi_m\in V'$. Prove
that the following three sets are equal to each other.
\begin{parts}
\item $\spn(\varphi_1,\dots,\varphi_m)$
\item $\bigl((\nul\varphi_1)\cap\cdots\cap(\nul\varphi_m)\bigr)^0$
\item $\{\varphi\in V' : (\nul\varphi_1)\cap\cdots\cap(\nul\varphi_m)
  \subseteq\nul\varphi\}$
\end{parts}
\end{exercise}
\begin{solution}
For ease of notation, let $A,B,C$ be the sets in the exercise, i.e.,
\begin{align*}
A &= \spn(\varphi_1,\dots,\varphi_m),\\
B &= \bigl((\nul\varphi_1)\cap\cdots\cap(\nul\varphi_m)\bigr)^0,\\
C &= \{\varphi\in V' \mid (\nul\varphi_1)\cap\cdots\cap(\nul\varphi_m)
      \subseteq\nul\varphi\}.
\end{align*}
We will show that $A=B$ and $B=C$. Starting off with $B=C$, note that by
definition,
\begin{align*}
B &= \bigl((\nul\varphi_1)\cap\cdots\cap(\nul\varphi_m)\bigr)^0\\
  &= \{\varphi\in V' \mid (\nul\varphi_1)\cap\cdots\cap(\nul\varphi_m)
     \subseteq\nul\varphi\} = C,
\end{align*}
where the reformulation of the annihilator subspace was shown in
Exercise~3F19.

To show that $A=B$, we may use Exercise~3F22(b) repeatedly to rewrite $B$ as
\[ B = (\nul\varphi_1)^0+\cdots+(\nul\varphi_m)^0. \]
By Exercise~3F6, we have $(\nul\varphi_i)^0=\spn(\varphi_i)$ for each
$i=1,\dots,m$. Therefore,
\begin{align*}
B &= (\nul\varphi_1)^0+\cdots+(\nul\varphi_m)^0\\
  &= \spn(\varphi_1)+\cdots+\spn(\varphi_m) = \spn(\varphi_1,\dots,\varphi_m)
   = A.
\end{align*}
We conclude that $A=B$.
\end{solution}

\begin{exercise}
Suppose $V$ is finite-dimensional and $v_1,\dots,v_m\in V$. Define a linear map
$\Gamma\colon V'\to\F^m$ by
$\Gamma(\varphi)=\bigl(\varphi(v_1),\dots,\varphi(v_m)\bigr)$.
\begin{parts}
\item Prove that $v_1,\dots,v_m$ spans $V$ if and only if $\Gamma$ is injective.
\item Prove that $v_1,\dots,v_m$ is linearly independent if and only if
  $\Gamma$ is surjective.
\end{parts}
\end{exercise}
\begin{solution}\leavevmode
\begin{parts}
\item Suppose $v_1,\dots,v_m$ spans $V$. If $\Gamma(\varphi)=0$, then
  $\varphi(v_i)=0$ for all $i=1,\dots,m$. Since $v_1,\dots,v_m$ spans $V$, it
  follows that $\varphi(v)=0$ for all $v\in V$, i.e., $\varphi=0$. Hence,
  $\Gamma$ is injective.

  Conversely, suppose $\Gamma$ is injective. If $v_1,\dots,v_m$ does not span
  $V$, then by Exercise~3F4, there exists a nonzero $\varphi\in V'$ that
  vanishes on $v_1,\dots,v_m$. But then $\Gamma(\varphi)=0$, contradicting the
  injectivity of $\Gamma$. Hence, $v_1,\dots,v_m$ must span $V$.
\item Suppose $v_1,\dots,v_m$ is linearly independent. Extend this list to a
  basis $v_1,\dots,v_m,\allowbreak v_{m+1},\dots,v_n$ of $V$. Define linear functionals
  $\varphi_1,\dots,\varphi_m\in V'$ by setting $\varphi_i(v_j)=\delta_{ij}$ for
  $i,j=1,\dots,m$ and $\varphi_i(v_j)=0$ for $j=m+1,\dots,n$. Then for any
  $(a_1,\dots,a_m)\in\F^m$, the linear functional
  $\varphi=a_1\varphi_1+\cdots+a_m\varphi_m$ satisfies
  $\Gamma(\varphi)=(a_1,\dots,a_m)$. Hence, $\Gamma$ is surjective.

  Conversely, suppose $\Gamma$ is surjective. If $v_1,\dots,v_m$ is not
  linearly independent, then there exist scalars $a_1,\dots,a_m$, not all zero,
  such that $a_1v_1+\cdots+a_mv_m=0$. Consider some index $j$ such that
  $a_j\ne0$. Since $\Gamma$ is surjective, there exists $\varphi\in V'$ such
  that $\varphi(v_j)=1$ and $\varphi(v_i)=0$ for all $i\ne j$. Then
  \[ 0 = \varphi(a_1v_1+\cdots+a_mv_m)
       = a_1\varphi(v_1)+\cdots+a_m\varphi(v_m) = a_j\ne0, \]
  a contradiction. Hence, $v_1,\dots,v_m$ must be linearly independent.
  \qedhere
\end{parts}
\end{solution}

\begin{exercise}
Suppose $V$ is finite-dimensional and $\varphi_1,\dots,\varphi_m\in V'$. Define
a linear map $\Gamma\colon V\to\F^m$ by
$\Gamma(v)=\bigl(\varphi_1(v),\dots,\varphi_m(v)\bigr)$.
\begin{parts}
\item Prove that $\varphi_1,\dots,\varphi_m$ spans $V'$ if and only if $\Gamma$
  is injective.
\item Prove that $\varphi_1,\dots,\varphi_m$ is linearly independent if and only
  if $\Gamma$ is surjective.
\end{parts}
\end{exercise}
\begin{solution}\leavevmode
\begin{parts}
\item Suppose $\varphi_1,\dots,\varphi_m$ spans $V'$. If $\Gamma(v)=0$, then
  $\varphi_i(v)=0$ for all $i=1,\dots,m$. Since $\varphi_1,\dots,\varphi_m$
  spans $V'$, then by Exercise~3F3, it follows that $\psi(v)=0$ for all
  $\psi\in V'$, i.e., $v=0$. Hence, $\Gamma$ is injective.

  Conversely, suppose $\Gamma$ is injective. Then $\nul\Gamma=\{0\}$. By
  definition, we have
  \[ \nul\Gamma = \nul\varphi_1\cap\cdots\cap\nul\varphi_m. \]
  By Exercise~3F23, we have
  \begin{align*}
  \spn(\varphi_1,\dots,\varphi_m)
   &= \bigl((\nul\varphi_1)\cap\cdots\cap(\nul\varphi_m)\bigr)^0\\
   &= (\nul\Gamma)^0 = \{\psi\in V' \mid \nul\Gamma\subseteq\nul\psi\} = V'.
  \end{align*}
  Hence, $\varphi_1,\dots,\varphi_m$ spans $V'$.
\item Suppose $\varphi_1,\dots,\varphi_m$ is linearly independent. Then
  $\dim\spn(\varphi_1,\dots,\varphi_m)=m$. By the fundamental theorem of linear
  maps and Exercise~3F23, we have
  \begin{align*}
  \dim\range\Gamma &= \dim V-\dim\nul\Gamma\\
   &= \dim V-\dim\bigl((\nul\varphi_1)\cap\cdots\cap(\nul\varphi_m)\bigr) = m.
  \end{align*}
  Since $\range\Gamma$ is a subspace of $\F^m$ and has dimension $m$, we
  conclude that $\range\Gamma=\F^m$, so $\Gamma$ is surjective.

  Conversely, suppose $\Gamma$ is surjective. Then $\range\Gamma=\F^m$, so
  $\dim\range\Gamma=m$. By the fundamental theorem of linear maps, we have
  \[ \dim\range\Gamma = \dim V-\dim\nul\Gamma
     = \dim V-\dim\bigl((\nul\varphi_1)\cap\cdots\cap(\nul\varphi_m)\bigr). \]
  By Exercise~3F23, we have
  \[ \spn(\varphi_1,\dots,\varphi_m)
     = \bigl((\nul\varphi_1)\cap\cdots\cap(\nul\varphi_m)\bigr)^0, \]
  so $\dim\spn(\varphi_1,\dots,\varphi_m)
  =\dim V-\dim\bigl((\nul\varphi_1)\cap\cdots\cap(\nul\varphi_m)\bigr)=m$.
  Hence, $\varphi_1,\dots,\varphi_m$ is linearly independent.
\end{parts}
\end{solution}

\begin{exercise}
Suppose $V$ is finite-dimensional and $\Omega$ is a subspace of $V'$. Prove that
\[ \Omega = \{v\in V : \varphi(v)=0 \text{ for every } \varphi\in\Omega\}^0. \]
\end{exercise}
\begin{solution}
\solnote{Manual Remark}{This says $\Omega=(\Omega^0)^0$.}

Let $S=\{v\in V \mid \varphi(v)=0 \text{ for every } \varphi\in\Omega\}$. We
need to show that $\Omega=S^0$. Let $\varphi_1,\dots,\varphi_m$ be a basis of
$\Omega$. Then $S=(\nul\varphi_1)\cap\cdots\cap(\nul\varphi_m)$. By
Exercise~3F23, we have
\[ \Omega = \spn(\varphi_1,\dots,\varphi_m)
   = \bigl((\nul\varphi_1)\cap\cdots\cap(\nul\varphi_m)\bigr)^0 = S^0.
   \qedhere \]
\end{solution}

\begin{exercise}
Suppose $T\in\Lin\bigl(\Poly_5(\R)\bigr)$ and $\nul T'=\spn(\varphi)$, where
$\varphi$ is the linear functional on $\Poly_5(\R)$ defined by
$\varphi(p)=p(8)$. Prove that
\[ \range T = \{p\in\Poly_5(\R) : p(8)=0\}. \]
\end{exercise}
\begin{solution}
Let $W=\{p\in\Poly_5(\R) \mid p(8)=0\}$. We need to show that $\range T=W$.

First, if $p\in\range T$, then there exists $q\in\Poly_5(\R)$ such that
$T(q)=p$. Since $\varphi\in\nul T'$, we have
\[ \varphi\bigl(T(q)\bigr) = \bigl(T'(\varphi)\bigr)(q) = 0. \]
But $\varphi\bigl(T(q)\bigr)=T(q)(8)=p(8)$, so $p(8)=0$. Hence, $p\in W$, and
we have shown that $\range T\subseteq W$.

Now we compare the dimensions of $\range T$ and $W$. Recall that
$(\range T)^0=\nul T'=\spn(\varphi)$, so $\dim(\range T)^0=1$. Since
$\Poly_5(\R)$ has dimension $6$, we have $\dim\range T=6-1=5$. On the other
hand, $W$ is the set of all polynomials in $\Poly_5(\R)$ that vanish at $8$.
Using the fundamental theorem of linear maps and the fact that $\varphi\ne0$ is
surjective by Exercise~3F1, we have $\dim W=\dim\Poly_5(\R)-1=5$. Hence,
$\range T$ and $W$ are subspaces of $\Poly_5(\R)$ with the same dimension, and
$\range T\subseteq W$. We conclude that $\range T=W$, as desired.
\end{solution}

\begin{exercise}
Suppose $V$ is finite-dimensional and $\varphi_1,\dots,\varphi_m$ is a linearly
independent list in $V'$. Prove that
\[ \dim\bigl((\nul\varphi_1)\cap\cdots\cap(\nul\varphi_m)\bigr)
   = (\dim V)-m. \]
\end{exercise}
\begin{solution}
Since $\varphi_1,\dots,\varphi_m$ is linearly independent, we have
$\dim\spn(\varphi_1,\allowbreak\dots,\varphi_m)=m$. By Exercise~3F23, we have
\[ \spn(\varphi_1,\dots,\varphi_m)
   = \bigl((\nul\varphi_1)\cap\cdots\cap(\nul\varphi_m)\bigr)^0. \]
Hence, $\dim\bigl((\nul\varphi_1)\cap\cdots\cap(\nul\varphi_m)\bigr)^0=m$.
Since $V$ is finite-dimensional, we have
\begin{align*}
\dim\bigl((\nul\varphi_1)\cap\cdots\cap(\nul\varphi_m)\bigr)
 &= \dim V-\dim\bigl((\nul\varphi_1)\cap\cdots\cap(\nul\varphi_m)\bigr)^0\\
 &= \dim V-m,
\end{align*}
as desired.
\end{solution}

\begin{exercise}
Suppose $V$ and $W$ are finite-dimensional and $T\in\Lin(V,W)$.
\begin{parts}
\item Prove that if $\varphi\in W'$ and $\nul T'=\spn(\varphi)$, then
  $\range T=\nul\varphi$.
\item Prove that if $\psi\in V'$ and $\range T'=\spn(\psi)$, then
  $\nul T=\nul\psi$.
\end{parts}
\end{exercise}
\begin{solution}\leavevmode
\begin{parts}
\item Since $(\range T)^0=\nul T'$ and $\nul T'=\spn(\varphi)$, we have
  $(\range T)^0=\spn(\varphi)$. Since $(U^0)^0=U$ in finite dimensions, taking
  annihilators yields
  \[ \range T = \bigl(\spn(\varphi)\bigr)^0 = \nul\varphi. \]
\item Using $\nul T=(\range T')^0$ and $\range T'=\spn(\psi)$ immediately
  yields
  \[ \nul T = \bigl(\spn(\psi)\bigr)^0 = \nul\psi. \qedhere \]
\end{parts}
\end{solution}

\begin{exercise}
Suppose $V$ is finite-dimensional and $\varphi_1,\dots,\varphi_n$ is a basis of
$V'$. Show that there exists a basis of $V$ whose dual basis is
$\varphi_1,\dots,\varphi_n$.
\end{exercise}
\begin{solution}
Let $\varphi_1,\dots,\varphi_n$ be a basis of $V'$. From Exercise~3F25, there
exists a map $\Gamma\colon V\to\F^n$ that is an isomorphism, since
$\varphi_1,\dots,\varphi_n$ is a basis for $V'$. Let $e_1,\dots,e_n$ be the
standard basis of $\F^n$. Since $\Gamma$ is an isomorphism, there exist
$v_1,\dots,v_n\in V$ such that $\Gamma(v_i)=e_i$ for each $i=1,\dots,n$. We
claim that $v_1,\dots,v_n$ is a basis of $V$ whose dual basis is
$\varphi_1,\dots,\varphi_n$.

To see that $v_1,\dots,v_n$ is a basis of $V$, we show linear independence.
Suppose $a_1,\dots,a_n$ are scalars such that $a_1v_1+\cdots+a_nv_n=0$.
Applying $\Gamma$ to both sides gives
\[ a_1e_1+\cdots+a_ne_n = \Gamma(a_1v_1+\cdots+a_nv_n) = \Gamma(0) = 0. \]
Since $e_1,\dots,e_n$ is a basis of $\F^n$, it follows that
$a_1=\cdots=a_n=0$. Hence, $v_1,\dots,v_n$ is linearly independent. Since
$\dim V=n$, we conclude that $v_1,\dots,v_n$ is a basis of $V$. To see that the
dual basis of $v_1,\dots,v_n$ is $\varphi_1,\dots,\varphi_n$, we need to show
that $\varphi_i(v_j)=\delta_{ij}$ for each $i,j=1,\dots,n$. Since
$\Gamma(v_j)=e_j$, we have
\[ \varphi_i(v_j) = \bigl(\Gamma(v_j)\bigr)_i = (e_j)_i = \delta_{ij}, \]
as desired.
\end{solution}

\begin{exercise}
Suppose $U$ is a subspace of $V$. Let $i\colon U\to V$ be the inclusion map
defined by $i(u)=u$. Thus $i'\in\Lin(V',U')$.
\begin{parts}
\item Show that $\nul i'=U^0$.
\item Prove that if $V$ is finite-dimensional, then $\range i'=U'$.
\item Prove that if $V$ is finite-dimensional, then $\tilde{\imath}'$ is an
  isomorphism from $V'/U^0$ onto $U'$.
\end{parts}
\exnote{The isomorphism in \textup{(c)} is natural in that it does not depend on
a choice of basis in either vector space.}
\end{exercise}
\begin{solution}\leavevmode
\begin{parts}
\item Note that if $\varphi\in\nul i'$, then $i'(\varphi)=0$, so
  $\varphi\bigl(i(u)\bigr)=0$ for all $u\in U$. Since $i(u)=u$, we have
  $\varphi(u)=0$ for all $u\in U$, so $\varphi\in U^0$. Hence,
  $\nul i'\subseteq U^0$. Conversely, if $\varphi\in U^0$, then $\varphi(u)=0$
  for all $u\in U$, so $i'(\varphi)=0$. Hence, $U^0\subseteq\nul i'$, and we
  conclude that $\nul i'=U^0$.
\item To see $\range i'\subseteq U'$, note that if $\varphi\in V'$, then
  $i'(\varphi)\in U'$ since $i'(\varphi)(u)=\varphi\bigl(i(u)\bigr)=\varphi(u)$
  for all $u\in U$. Hence, $\range i'\subseteq U'$.

  Since $V$ is finite-dimensional, we have $\dim V'=\dim V$. By the fundamental
  theorem of linear maps, we have
  \[ \dim\range i' = \dim V'-\dim\nul i' = \dim V-\dim U^0. \]
  By the annihilator dimension formula, we have $\dim U^0=\dim V-\dim U$.
  Hence,
  \[ \dim\range i' = \dim V-(\dim V-\dim U) = \dim U = \dim U'. \]
  Since $\range i'$ is a subspace of $U'$ and has the same dimension, we
  conclude that $\range i'=U'$, as desired.
\item From part (a), we have $\nul i'=U^0$. Then we know that $V'/U^0$ is
  isomorphic to $\range i'$. Then the induced map
  $\tilde{\imath}'\colon V'/U^0\to\range i'$ defined by
  $\tilde{\imath}'(\varphi+U^0)=i'(\varphi)$ is an isomorphism. Since
  $\range i'=U'$ by part (b), we conclude that $i'$ induces an isomorphism from
  $V'/U^0$ onto $U'$, as desired.
\end{parts}
\end{solution}

\begin{exercise}
The \emph{double dual space} of $V$, denoted by $V''$, is defined to be the
dual space of $V'$. In other words, $V''=(V')'$. Define $\Lambda\colon V\to V''$
by
\[ (\Lambda v)(\varphi) = \varphi(v) \]
for each $v\in V$ and each $\varphi\in V'$.
\begin{parts}
\item Show that $\Lambda$ is a linear map from $V$ to $V''$.
\item Show that if $T\in\Lin(V)$, then $T''\circ\Lambda=\Lambda\circ T$, where
  $T''=(T')'$.
\item Show that if $V$ is finite-dimensional, then $\Lambda$ is an isomorphism
  from $V$ onto $V''$.
\end{parts}
\exnote{Suppose $V$ is finite-dimensional. Then $V$ and $V'$ are isomorphic, but
finding an isomorphism from $V$ onto $V'$ generally requires choosing a basis of
$V$. In contrast, the isomorphism $\Lambda$ from $V$ onto $V''$ does not require
a choice of basis and thus is considered more natural.}
\end{exercise}
\begin{solution}\leavevmode
\begin{parts}
\item To see that $\Lambda$ is linear, let $v,w\in V$ and $a\in\F$. Then for
  any $\varphi\in V'$, we have
  \begin{align*}
  \bigl(\Lambda(v+w)\bigr)(\varphi) &= \varphi(v+w) = \varphi(v)+\varphi(w)
    = (\Lambda v)(\varphi)+(\Lambda w)(\varphi),\\
  \bigl(\Lambda(av)\bigr)(\varphi) &= \varphi(av) = a\varphi(v)
    = a(\Lambda v)(\varphi).
  \end{align*}
  Hence, $\Lambda$ is linear.
\item Let $v\in V$ and $\varphi\in V'$. Then
  \begin{align*}
  \bigl((T''\circ\Lambda)(v)\bigr)(\varphi)
   &= \bigl(T''(\Lambda v)\bigr)(\varphi)\\
   &= (\Lambda v)\bigl(T'(\varphi)\bigr)\\
   &= T'(\varphi)(v)\\
   &= \varphi\bigl(T(v)\bigr)\\
   &= \Bigl(\Lambda\bigl(T(v)\bigr)\Bigr)(\varphi)
    = \bigl((\Lambda\circ T)(v)\bigr)(\varphi).
  \end{align*}
  Hence, $T''\circ\Lambda=\Lambda\circ T$.
\item Since $V$ is finite-dimensional, we have $\dim V=\dim V'$. Then
  $\dim V''=\dim(V')'=\dim V'=\dim V$. By the fundamental theorem of linear
  maps, we have
  \[ \dim\range\Lambda = \dim V-\dim\nul\Lambda. \]
  To see that $\Lambda$ is injective, suppose $\Lambda v=0$ for some $v\in V$.
  Then for any $\varphi\in V'$, we have $\varphi(v)=(\Lambda v)(\varphi)=0$.
  Since this holds for all $\varphi\in V'$ by Exercise~3F3, it follows that
  $v=0$. Hence, $\Lambda$ is injective, so $\nul\Lambda=\{0\}$. Therefore,
  $\dim\range\Lambda=\dim V-0=\dim V$. Since $\range\Lambda$ is a subspace of
  $V''$ and has the same dimension as $V''$, we conclude that
  $\range\Lambda=V''$, so $\Lambda$ is surjective. Combining injectivity and
  surjectivity, we conclude that $\Lambda$ is an isomorphism from $V$ onto
  $V''$, as desired.
\end{parts}
\end{solution}

\begin{exercise}
Suppose $U$ is a subspace of $V$. Let $\pi\colon V\to V/U$ be the usual quotient
map. Thus $\pi'\in\Lin\bigl((V/U)',V'\bigr)$.
\begin{parts}
\item Show that $\pi'$ is injective.
\item Show that $\range\pi'=U^0$.
\item Conclude that $\pi'$ is an isomorphism from $(V/U)'$ onto $U^0$.
\end{parts}
\exnote{The isomorphism in \textup{(c)} is natural in that it does not depend on
a choice of basis in either vector space. In fact, there is no assumption here
that any of these vector spaces are finite-dimensional.}
\end{exercise}
\begin{solution}\leavevmode
\begin{parts}
\item Suppose $\pi'(\varphi)=0$ for some $\varphi\in(V/U)'$. Then
  $\varphi\bigl(\pi(v)\bigr)=0$ for all $v\in V$. Since $\pi(v)=v+U$, we have
  $\varphi(v+U)=0$ for all $v\in V$. Since this holds for all $v\in V$ by
  surjectivity of $\pi$, it follows that $\varphi=0$. Hence, $\pi'$ is
  injective.
\item To see $\range\pi'\subseteq U^0$, note that if $\varphi\in(V/U)'$, then
  for any $u\in U$, we have
  $\pi'(\varphi)(u)=\varphi\bigl(\pi(u)\bigr)=\varphi(0+U)=0$. Hence,
  $\range\pi'\subseteq U^0$.

  To see the other inclusion, let $\psi\in U^0$. We want to find
  $\varphi\in(V/U)'$ such that $\pi'(\varphi)=\psi$. Define
  $\varphi\colon V/U\to\F$ by $\varphi(v+U)=\psi(v)$ for each $v\in V$. We need
  to check that $\varphi$ is well-defined. If $v+U=w+U$, then $v-w\in U$, so
  $\psi(v-w)=0$, which implies that $\psi(v)=\psi(w)$. Hence, $\varphi$ is
  well-defined. It is also linear since for any $a,b\in\F$ and $v,w\in V$, we
  have
  \begin{align*}
  \varphi\bigl(a(v+U)+b(w+U)\bigr) &= \varphi\bigl((av+bw)+U\bigr)
    = \psi(av+bw)\\
   &= a\psi(v)+b\psi(w) = a\varphi(v+U)+b\varphi(w+U).
  \end{align*}
  Finally, we have
  $\pi'(\varphi)(v)=\varphi\bigl(\pi(v)\bigr)=\varphi(v+U)=\psi(v)$ for all
  $v\in V$, so $\pi'(\varphi)=\psi$. Hence, $U^0\subseteq\range\pi'$, and we
  conclude that $\range\pi'=U^0$.
\item From part (a), we have $\pi'$ is injective. From part (b), we have
  $\range\pi'=U^0$. Then the induced map $\tilde{\pi}'\colon(V/U)'\to U^0$
  defined by $\tilde{\pi}'(\varphi)=\pi'(\varphi)$ is an isomorphism. Hence,
  $\pi'$ is an isomorphism from $(V/U)'$ onto $U^0$, as desired.
\end{parts}
\end{solution}
